\begin{textsample}{POS Dim 3 – rs – Score 22.00 – rs/rs\_inf\_10.txt}  \label{ex:f3_pos_221}
10.
DIREITOS DE APRENDIZAGEM E DESENVOLVIMENTO E CAMPOS DE EXPERIÊNCIAS NO REFERENCIAL CURRICULAR GAÚCHO Na Educação \textbf{Infantil}, as interações e a \textbf{brincadeira} compõem o eixo estruturante das propostas pedagógicas, assegurando às \textbf{crianças} os direitos de conviver, \textbf{brincar}, participar, explorar, expressar e conhecer -se.
Os seis Direitos de Aprendizagem e Desenvolvimento, de acordo com a BNCC, > (... ) asseguram as condições para que as \textbf{crianças} aprendam em situações nas quais possam desempenhar um papel ativo em ambientes que as convidem a vivenciar desafios e a sentirem -se provocadas a resolvê -los, nas quais possam construir significados sobre si, os outros e o mundo social e natural. ( BNCC, 2017, p. 35 ).
Os Direitos de Aprendizagem e Desenvolvimento fundamentam -se nos princípios éticos, políticos e estéticos estabelecidos pelas DCNEI, que orientam as propostas pedagógicas das instituições de Educação \textbf{Infantil} de todo o país.
São eles : CONVIVER com outras \textbf{crianças} e \textbf{adultos}, em \textbf{pequenos} e grandes grupos, utilizando diferentes linguagens, ampliando o conhecimento de si e do outro, o respeito em relação à cultura e às diferenças entre as pessoas. \textbf{BRINCAR} de diversas formas, em diferentes espaços e tempos, com diferentes parceiros ( \textbf{crianças} e \textbf{adultos} ) de forma a ampliar e diversificar suas possibilidades de acesso a produções culturais.
A participação e as transformações introduzidas pelas \textbf{crianças} nas \textbf{brincadeiras} devem ser valorizadas, tendo em vista o estímulo ao desenvolvimento de seus conhecimentos, sua imaginação, criatividade, experiências emocionais, corporais, \textbf{sensoriais}, expressivas, cognitivas, sociais e relacionais.
PARTICIPAR ativamente, com \textbf{adultos} e outras \textbf{crianças}, tanto do planejamento da gestão da escola e das atividades propostas pelo educador quanto da realização das atividades da vida cotidiana, tais como a escolha das \textbf{brincadeiras}, dos materiais e dos ambientes, desenvolvendo diferentes linguagens e elaborando conhecimentos, decidindo e se posicionando.
EXPLORAR movimentos, \textbf{gestos}, \textbf{sons}, formas, \textbf{texturas}, cores, palavras, emoções, transformações, relacionamentos, histórias, objetos, elementos da natureza, na escola e fora dela, ampliando seus saberes sobre a cultura em suas diversas modalidades : as artes, a escrita, a ciência e a tecnologia.
EXPRESSAR, como sujeito dialógico, criativo e sensível, suas necessidades, emoções, sentimentos, dúvidas, hipóteses, descobertas, opiniões, questionamentos, por meio de diferentes linguagens.
CONHECER -SE e construir sua identidade pessoal, social e cultural, constituindo uma imagem positiva de si e de seus grupos de pertencimento, nas diversas experiências de \textbf{cuidados}, interações, \textbf{brincadeiras} e linguagens vivenciadas na instituição escolar e em seu contexto familiar e comunitário.
A figura abaixo ilustra a inter-relação entre os seis Direitos de Aprendizagem e Desenvolvimento e os princípios orientadores da Educação \textbf{Infantil} : A concepção de uma \textbf{criança} ativa e capaz, com direitos e \textbf{desejos}, sujeito que observa, questiona, levanta hipóteses, constrói e se apropria de conhecimentos por meio da ação e nas interações com o mundo físico e social convoca a refletir e a modificar as formas tradicionais de planejar e desenvolver as práticas pedagógicas, diferenciando -se da estrutura baseada em áreas do conhecimento, mais familiar aos currículos efetivados no Ensino Fundamental e Médio.
A BNCC, de modo a orientar os Projetos Político-Pedagógicos ( PPPs ) das instituições de Educação \textbf{Infantil}, propõe que nos Campos de Experiências as \textbf{crianças} tenham garantidos os seis Direitos de Aprendizagem e Desenvolvimento mediadores de significativas aprendizagens.
Estes Direitos são retomados em cada Campo de Experiências e são a referência para a elaboração de objetivos de aprendizagem e desenvolvimento e para o planejamento do professor.
A estruturação curricular da Educação \textbf{Infantil} está organizada em cinco Campos de Experiências, conforme proposto na BNCC ( 2017 ).
Os Campos de Experiências constituem um arranjo curricular que \textbf{acolhe} as situações e experiências concretas de vida das \textbf{crianças} e seus saberes, os diversos contextos das culturas locais e regionais e articula -os aos conhecimentos que fazem parte do patrimônio que a humanidade produziu.
Na ideia dos Campos de Experiências, reside a articulação das dimensões do conhecimento, das práticas sociais e das múltiplas linguagens.
Na figura abaixo, evidencia -se a relação entre os diferentes Campos de Experiências : Assim, a organização curricular por Campos de Experiências é fundamentada em uma concepção de \textbf{criança} que age, cria, produz sentidos sobre si e sobre o mundo e aprende nas relações e experiências que vive, de maneira integrada.
Portanto, os Campos de Experiências subvertem a lógica disciplinar de estruturar o conhecimento, centrando -se na produção de saberes das \textbf{crianças} que são sustentados pelas relações e interações.
Daí a importância de práticas educativas que valorizam experiências concretas da vida cotidiana.
Desse modo, a proposta curricular do Estado do Rio Grande do Sul para a Educação \textbf{Infantil} organiza -se conforme o estabelecido na BNCC ( 2017 ) e nas DCNEI ( 2009 ), em seu Art. 3º, onde destaca que : > O currículo da Educação \textbf{Infantil} é concebido como um conjunto de práticas que buscam articular as experiências e os saberes das \textbf{crianças} com os conhecimentos que fazem parte do patrimônio cultural, artístico, ambiental, científico e tecnológico, de modo a promover o desenvolvimento integral de \textbf{crianças} de 0 a 5 anos de idade. ( BRASIL, 2009 ).
A partir dessas recomendações, é preciso considerar que a aprendizagem tem, como ponto de partida, o que a \textbf{criança} já sabe e o que ela é capaz de fazer.
O conceito de experiência reconhece que a imersão das \textbf{crianças} em práticas sociais e culturais criativas e interativas promove significativas aprendizagens criando momentos plenos de afetividade e descobertas.
A presença de um professor sensível e atento é fundamental para que as \textbf{crianças} vivam experiências mediadoras de valiosas aprendizagens em que expressem seus \textbf{desejos} e descobertas pelo corpo, \textbf{gestos} e palavras.
Cabe ao professor proporcionar experiências ricas, desafiadoras e variadas que possibilitem a cada \textbf{criança} desenvolver seu próprio percurso educativo, que é único e fruto de uma variedade de experiências que as \textbf{crianças} vivenciam na escola.
Isso significa que cada \textbf{criança} tem um potencial de desenvolvimento sobre o qual o professor deve atuar.
A organização curricular por Campos de Experiências propõe que as ações pedagógicas sejam desenvolvidas a partir de uma \textbf{escuta} atenta sobre as \textbf{crianças}, colocando em relação aos saberes das \textbf{crianças} e os saberes dos professores, por meio de uma pedagogia relacional, em que o conhecimento é construído na interação entre as \textbf{crianças}, com os \textbf{adultos} e com o mundo.
Nesse sentido, as ações planejadas pelo professor devem ser marcadas pela intencionalidade educativa na organização de experiências que permitam às \textbf{crianças} articular e conhecer a si, o outro, a natureza, a cultura e a produção científica por meio das interações e da \textbf{brincadeira} ( BRASIL, 2017 ).
Como orienta a BNCC : > Parte do trabalho do educador é refletir, selecionar, organizar, planejar, mediar e monitorar o conjunto das práticas e interações, garantindo a pluralidade de situações que promovam o desenvolvimento pleno das \textbf{crianças}. ( BRASIL, 2017, p. 37 ).
Pensar um currículo organizado por Campos de Experiências é compreender que esses Campos articulam -se entre si, que não há uma fragmentação ou divisão disciplinar entre os Campos ; é reconhecer que as \textbf{crianças} têm em si o \textbf{desejo} de aprender e que o papel do \textbf{adulto} passa por desconstruir algumas práticas tradicionais e enrijecidas e construir novas práticas que possibilitem às \textbf{crianças} dar significado aos diferentes contextos de interação e que possam representar, em suas \textbf{brincadeiras}, diferentes fatos de suas vivências.
Em outras palavras, os Campos de Experiências expressam a forma interdisciplinar que os conhecimentos são produzidos.
Os Campos de Experiências podem subsidiar as práticas das \textbf{crianças} isoladamente ou reunindo os objetivos de um ou mais Campos, e envolvem todos os momentos da jornada das \textbf{crianças} na Educação \textbf{Infantil}, incluindo o acolhimento inicial, o momento das \textbf{refeições}, a participação delas no planejamento das atividades, as festividades e encontros com as famílias, as atividades de expressão, investigação, as \textbf{brincadeiras}, realizadas ao longo da jornada \textbf{diária} e semanal das \textbf{crianças}.
Dessa maneira, os referidos Campos não são trabalhados apenas em um dia marcado da semana, nem há expectativa de haver uma aula de 45 minutos para o trabalho com um Campo em cada dia ou para que determinado bimestre do ano letivo seja dedicado apenas a um Campo.
É importante destacar que as experiências que perpassam pelos diversos Campos consideram que as \textbf{crianças} estão descobrindo como é estar no mundo, como as coisas funcionam e como podem ser chamadas.
Por isso, as práticas sociais e da cultura são aprendizagens que ganham significado e compõem os Campos de Experiências a serem contemplados na organização curricular da escola da infância : as acolhidas e transições \textbf{diárias}, a alimentação, a \textbf{higiene}, o repouso, o convívio com outras \textbf{crianças} e \textbf{adultos}, as \textbf{brincadeiras} e a ampliação de repertórios da cultura por meio da articulação de saberes das \textbf{crianças} com os conhecimentos que a humanidade já sistematizou.
Em conformidade com a BNCC ( 2017 ), são cinco os Campos de Experiências para os \textbf{bebês}, as \textbf{crianças} bem \textbf{pequenas} e as \textbf{crianças} \textbf{pequenas} : - O Eu, o Outro e o Nós ; - Corpo, \textbf{Gestos} e Movimentos ; - \textbf{Escuta}, Fala, Pensamento e Imaginação ; - Traços, \textbf{Sons}, Cores e Formas ; - Espaços, Tempos, Quantidades, Relações e Transformações.
Além de assegurados os Direitos de Aprendizagem e Desenvolvimento e a organização curricular por Campos de Experiências, reconhece -se que cada faixa etária possui especificidades quanto às possibilidades de aprendizagem e às características do desenvolvimento das \textbf{crianças}.
É importante destacar que os grupos etários estabelecidos na BNCC ( 2017 ) não devem ser considerados de forma rígida, pois as \textbf{crianças} se desenvolvem e aprendem de acordo com ritmos próprios que precisam ser observados nas práticas pedagógicas.
Assim, a BNCC define três grupos etários, a partir dos quais constituem -se os Objetivos de aprendizagem e desenvolvimento.
São eles : | \textbf{CRECHE} — \textbf{Bebês} | \textbf{CRECHE} — \textbf{Crianças} bem \textbf{pequenas} | PRÉ-ESCOLA — \textbf{Crianças} \textbf{pequenas} | |---|---|---| | 0 a 1 ano e 6 \textbf{meses} | 1 ano e 7 \textbf{meses} a 3 anos e 11 \textbf{meses} | 4 anos a 5 anos e 11 \textbf{meses} | FONTE : Base Nacional Comum Curricular, 2017, p. 42 Os Objetivos de aprendizagem e desenvolvimento propostos para cada grupo etário são apresentados nos Campos de Experiências, como subsídios para o planejamento das práticas pedagógicas.
Os parâmetros para a organização dos grupos têm como referência a faixa etária e a proposta pedagógica das instituições, observada a legislação vigente.
O Referencial Curricular Gaúcho da Educação \textbf{Infantil} organiza -se de acordo com a BNCC, em que cada objetivo de aprendizagem e desenvolvimento aparece identificado por um código alfanumérico, acrescido do código do objetivo do território gaúcho, com a seguinte composição, como \textbf{demonstra} a figura a seguir : O esquema acima \textbf{demonstra} como os Objetivos de aprendizagem e desenvolvimento são indicados no documento. - As duas primeiras letras ( EI ) indicam a primeira etapa da Educação Básica, a Educação \textbf{Infantil}. - Os dois primeiros números indicam o grupo por faixa etária, ou seja, 01 = \textbf{Bebês} ( zero a 1 ano e 6 \textbf{meses} ), 02 = \textbf{Crianças} bem \textbf{pequenas} ( 1 ano e 7 \textbf{meses} a 3 anos e 11 \textbf{meses} ) e 03 = \textbf{Crianças} \textbf{pequenas} ( 4 anos a 5 anos e 11 \textbf{meses} ). - O segundo par de letras indica um dos Campos de Experiências : EO = O Eu, o Outro e o Nós ; CG = Corpo, \textbf{Gestos} e Movimentos ; TS = Traços, \textbf{Sons}, Cores e Formas ; EF = \textbf{Escuta}, Fala, Pensamento e Imaginação ; ET = Espaços, Tempos, Quantidades, Relações e Transformações. - Os dois números seguintes indicam a posição do Objetivo na numeração sequencial do Campo de Experiências para cada grupo etário ; no entanto a sequência dos códigos alfanuméricos não sugere ordem ou hierarquia entre os objetivos de aprendizagem e desenvolvimento. - O terceiro par de letras ( RS ) indica o Estado do Rio Grande do Sul. - Os dois últimos números indicam a posição do objetivo na numeração do Campo de Experiências para cada grupo/faixa etária dentro do território gaúcho.
É importante destacar que os Campos de Experiências não são lineares, ou seja, não obedecem a uma ordem de prioridades, mas articulam -se entre si.
Ao realizar a leitura dos Objetivos de aprendizagem e desenvolvimento na Educação \textbf{Infantil}, o professor deve atentar para o contínuo das aprendizagens no grupo etário ( progressão vertical ) e entre os grupos etários ( progressão horizontal ), preocupando -se com a inter-relação entre os campos, evitando a fragmentação e a descontinuidade do trabalho pedagógico.
A seguir, são apresentados cada um dos Campos de Experiências e seus respectivos Objetivos de aprendizagem e desenvolvimento, elaborados mediante a participação de professores e profissionais da Educação \textbf{Infantil} de todo o Estado do Rio Grande do Sul, levando -se em conta as especificidades e características dos contextos regionais e locais, de forma a contribuir para a potencialização das práticas pedagógicas nas instituições de Educação \textbf{Infantil} do território gaúcho.

% matched lemmas: acolher, adulto, bebê, brincadeira, brincar, creche, criança, cuidado, demonstrar, desejo, diário, escuta, gesto, higiene, infantil, mês, pequeno, refeição, sensorial, som, textura
\end{textsample}
